% Fig — the five moves that collapse the PDE to the closure ODE.
% Mirrors app:pde2ode. Body only.
\begin{tikzpicture}[font=\footnotesize, every node/.style={align=center},
  st/.style={rectangle, rounded corners=2pt, draw=dim, fill=white,
             line width=0.8pt, inner sep=5pt, text width=118mm},
  res/.style={st, draw=forest, fill=forest!10, line width=1pt},
  arr/.style={-{Stealth[length=2.2mm]}, line width=0.9pt, draw=dim, line cap=round}]
  \node[st] (m0) at (0,0)
    {\textbf{the PDE} \quad $\rho c_p\,\partial_t T=\partial_z(K\,\partial_z T)$
     \hfill {\scriptsize\color{dim}two variables: $z$ and $t$}};
  \node[st] (m1) at (0,-1.32)
    {\textbf{1\; average over one lunation} \quad apply
     $\langle\cdot\rangle=\frac1P\int_0^P dt$ to both sides};
  \node[st] (m2) at (0,-2.62)
    {\textbf{2\; the time term dies} \quad
     $\langle\partial_t h\rangle=\frac1P[h(P)-h(0)]=0$ --- the cycle closes\\[1pt]
     {\scriptsize\color{dim}legal ONLY in periodic steady state}};
  \node[st] (m3) at (0,-3.98)
    {\textbf{3\; average commutes with} $\partial_z$ \quad
     $0=\partial_z\langle K\,\partial_z T\rangle$};
  \node[st] (m4) at (0,-5.28)
    {\textbf{4\; pin the constant at the base} \quad
     $\langle K\,\partial_z T\rangle=Q_b$ at every depth};
  \node[res] (m5) at (0,-6.85)
    {\textbf{5\; peel the average} (name the gap $u_{\rm rect}$):\\[2pt]
     $\dfrac{d\Tm}{dz}=\dfrac{Q_b-u_{\rm rect}}{K(\Tm,z)}$\\[2pt]
     {\scriptsize one variable: an ODE in depth --- Step B}};
  \foreach \a/\b in {m0/m1, m1/m2, m2/m3, m3/m4, m4/m5} \draw[arr] (\a) -- (\b);
\end{tikzpicture}
